% ============================================================================
% APÉNDICE B: BIBLIOGRAFÍA Y FUENTES
% ============================================================================
\chapter{Bibliografía y Fuentes}
\label{ap:bibliografia}

\section{Ondas Gravitacionales --- LIGO/Virgo}

\begin{enumerate}
    \item Abbott, B. P., et al. (LIGO/Virgo). ``Observation of Gravitational Waves from a Binary Black Hole Merger.'' Physical Review Letters 116, 061102 (2016). arXiv: 1602.03837

    \item Abbott, B. P., et al. ``Tests of General Relativity with GW150914.'' Physical Review Letters 116, 221101 (2016). arXiv: 1602.03841

    \item Abedi, J., Dykaar, H., \& Afshordi, N. ``Echoes from the Abyss: Tentative evidence for Planck-scale structure at black hole horizons.'' Physical Review D 96, 082004 (2017). arXiv: 1612.00266
\end{enumerate}


\section{Topología --- Botella de Klein}

\begin{enumerate}
    \item Klein, F. ``Über Riemann's Theorie der algebraischen Funktionen und ihrer Integrale.'' (1882). Leipzig: B.G. Teubner. [Trabajo original]

    \item Nakahara, M. ``Geometry, Topology and Physics.'' 2nd Edition. CRC Press (2003). ISBN: 978-0750306065

    \item Hatcher, A. ``Algebraic Topology.'' Cambridge University Press (2002). Disponible: \url{https://pi.math.cornell.edu/~hatcher/AT/ATpage.html}
\end{enumerate}


\section{Dimensiones Extra --- Kaluza-Klein}

\begin{enumerate}
    \item Kaluza, T. ``Zum Unitätsproblem der Physik.'' Sitzungsberichte der Preussischen Akademie der Wissenschaften, 966-972 (1921).

    \item Klein, O. ``Quantentheorie und fünfdimensionale Relativitätstheorie.'' Zeitschrift für Physik 37, 895-906 (1926).

    \item Randall, L., \& Sundrum, R. ``A Large Mass Hierarchy from a Small Extra Dimension.'' Physical Review Letters 83, 3370 (1999). arXiv: hep-ph/9905221

    \item Polchinski, J. ``String Theory.'' Cambridge University Press (1998). Vol. I \& II.
\end{enumerate}


\section{Constantes Fundamentales --- Fuentes de Datos}

\begin{enumerate}
    \item Tiesinga, E., et al. ``CODATA Recommended Values of the Fundamental Physical Constants: 2018.'' Reviews of Modern Physics 93, 025010 (2021). \url{https://physics.nist.gov/cuu/Constants/}

    \item Particle Data Group. ``Review of Particle Physics.'' Progress of Theoretical and Experimental Physics 2022, 083C01 (2022). \url{https://pdg.lbl.gov/}
\end{enumerate}


\section{Cosmología --- CMB y Constante Cosmológica}

\begin{enumerate}
    \item Planck Collaboration. ``Planck 2018 results. VI. Cosmological parameters.'' Astronomy \& Astrophysics 641, A6 (2020). arXiv: 1807.06209

    \item Penzias, A. A., \& Wilson, R. W. ``A Measurement of Excess Antenna Temperature at 4080 Mc/s.'' Astrophysical Journal 142, 419-421 (1965). [Premio Nobel 1978]

    \item Weinberg, S. ``The Cosmological Constant Problem.'' Reviews of Modern Physics 61, 1-23 (1989).

    \item Riess, A. G., et al. ``Observational Evidence from Supernovae for an Accelerating Universe.'' Astronomical Journal 116, 1009-1038 (1998). [Premio Nobel 2011]
\end{enumerate}


\section{Física de Partículas --- Modelo Estándar}

\begin{enumerate}
    \item Weinberg, S. ``A Model of Leptons.'' Physical Review Letters 19, 1264 (1967). [Premio Nobel 1979]

    \item Higgs, P. W. ``Broken Symmetries and the Masses of Gauge Bosons.'' Physical Review Letters 13, 508 (1964). [Premio Nobel 2013]

    \item ATLAS/CMS Collaborations. ``Observation of a new particle at 125 GeV.'' Physics Letters B 716 (2012).

    \item Georgi, H., \& Glashow, S. L. ``Unity of All Elementary-Particle Forces.'' Physical Review Letters 32, 438 (1974). [SU(5) GUT]
\end{enumerate}


\section{Neutrinos}

\begin{enumerate}
    \item Fukuda, Y., et al. (Super-Kamiokande). ``Evidence for Oscillation of Atmospheric Neutrinos.'' Physical Review Letters 81, 1562 (1998). [Premio Nobel 2015]

    \item An, F. P., et al. (Daya Bay). ``Observation of Electron-Antineutrino Disappearance.'' Physical Review Letters 108, 171803 (2012). [$\theta_{13}$]

    \item NuFIT Collaboration. \url{http://www.nu-fit.org/}
\end{enumerate}


\section{Asimetría Materia-Antimateria}

\begin{enumerate}
    \item Sakharov, A. D. ``Violation of CP Invariance, C Asymmetry, and Baryon Asymmetry of the Universe.'' JETP Letters 5, 24-27 (1967). [Condiciones de Sakharov]

    \item Christenson, J. H., et al. ``Evidence for the $2\pi$ Decay of the $K_2^0$ Meson.'' Physical Review Letters 13, 138 (1964). [Premio Nobel 1980]
\end{enumerate}


\section{Relatividad General y Electromagnetismo}

\begin{enumerate}
    \item Einstein, A. ``Die Grundlage der allgemeinen Relativitätstheorie.'' Annalen der Physik 354, 769-822 (1916).

    \item Misner, C. W., Thorne, K. S., \& Wheeler, J. A. ``Gravitation.'' W. H. Freeman (1973).

    \item Maxwell, J. C. ``A Dynamical Theory of the Electromagnetic Field.'' Phil. Trans. Royal Soc. 155, 459-512 (1865).

    \item Jackson, J. D. ``Classical Electrodynamics.'' 3rd Edition. Wiley (1999).
\end{enumerate}


\section{Valores Experimentales Utilizados}

\begin{table}[H]
\centering
\begin{tabular}{ll}
\toprule
\textbf{Constante} & \textbf{Valor} \\
\midrule
Velocidad de la luz & $c = 299,792,458$ m/s (exacto) \\
Estructura fina & $1/\alpha = 137.035999084(21)$ \\
Masa del electrón & $m_e = 9.1093837015(28) \times 10^{-31}$ kg \\
Masa del protón & $m_p = 1.67262192369(51) \times 10^{-27}$ kg \\
Masa del Higgs & $m_H = 125.25 \pm 0.17$ GeV/$c^2$ \\
Temperatura CMB & $T_{\text{CMB}} = 2.7255 \pm 0.0006$ K \\
Asimetría bariónica & $\eta_B = (6.1 \pm 0.2) \times 10^{-10}$ \\
Número de Avogadro & $N_A = 6.02214076 \times 10^{23}$ mol$^{-1}$ (exacto) \\
\bottomrule
\end{tabular}
\end{table}


\section{Nota sobre Originalidad}

La Teoría Klein presentada en este libro es \textbf{trabajo original} de Fausto José Di Bacco. Las referencias aquí listadas proporcionan contexto histórico, valores experimentales y fundamentos teóricos establecidos.

Las \textbf{fórmulas nuevas} ($7\pi$, $6\pi^5$, $7^2\pi - 7 - \pi^2$, etc.) son \textbf{contribuciones originales} que no aparecen en la literatura previa.
